\chapter{Описание на реализацията}

\section{Разпределение на типовете по слоеве}

За да бъде архитектурата ясна и лесна за поддръжка, типовете в проекта са разпределени по слоеве според тяхната функция. В домейн слоя се намират основните обектни типове, които описват същинските бизнес обекти: видео, анализ и резултат от детекция. Тези класове не познават конкретна технология за потребителски интерфейс, нито начин на съхранение, което позволява да служат като стабилен модел на системата.

Приложният слой съдържа класовете, които оформят интерфейса между домейна и технологията за представяне. Тук се намират DTO-тата, моделите на заявки, service contract-ите и споделените ViewModel-и като \texttt{HomeViewModel}, \texttt{VideoListViewModel}, \texttt{CreateVideoViewModel}, \texttt{AnalysisRunListViewModel}, \texttt{CreateAnalysisRunViewModel} и \texttt{AnalysisRunDetailsViewModel}. В същия слой са разположени и преизползваемите абстракции за филтри и за списъци от обекти, които позволяват една и съща логика да се прилага върху различни екрани.

Infrastructure слоят съдържа реалните имплементации, върху които стъпва приложението: \texttt{AppDbContext}, EF Core service класовете, услугите за съхранение на файлове, логиката за обработка на видео, фоновата услуга и модулите за инференция. Така проектът остава ясно структуриран: приложният слой описва какво трябва да може системата, а инфраструктурният слой реализира как точно става това.

\section{Организация на преизползването на ViewModel-и}

Изискването за едни и същи ViewModel-и в две различни View технологии е едно от най-съществените за проекта. Вместо две паралелни дървета от класове със сходна функционалност, едното за уеб, а другото за настолен интерфейс, споделените ViewModel-и се намират в \texttt{CVAnalysisHub.Application} и се използват и от двата клиента.

Това преизползване е приложено в основните функционални сценарии на приложението. Обзорните метрики, списъците от видеа и анализи, филтрите, процесът на качване, процесът по създаване на анализ и детайлният преглед на избран analysis run са реализирани чрез споделени приложни ViewModel-и. Avalonia има собствен \texttt{MainWindowViewModel}, но неговата роля не е да дублира бизнес функционалността, а да координира специфичното за настолния клиент състояние на интерфейса върху тези споделени ViewModel-и, например предварителни кадри и действия за отваряне на изходното и обработеното видео.

\section{Организация на преизползването за двете СУБД}

Поддръжката на две СУБД е реализирана така, че да не се нарушава идеята за общ модел и обща бизнес логика. Не съществува отделен service слой за SQLite и отделен service слой за PostgreSQL. Вместо това изборът на доставчик се извършва конфигурационно в инфраструктурния слой. Там, въз основа на стойността на \texttt{Persistence:Provider}, се избира дали EF Core да използва \texttt{UseSqlite} или \texttt{UseNpgsql}.

Тази организация избягва дублиране на една и съща логика за всяка СУБД. В текущия проект домейн и приложният слой не се интересуват кой точно доставчик стои отдолу. За тях услугата за съхранение е една и съща по поведение, което позволява системата да бъде стартирана и с SQLite, и с PostgreSQL без промяна във ViewModel-ите и без пренаписване на бизнес логиката.

\section{Организация на преизползваемите UI механизми}

Преизползваемата функционалност за филтриране е реализирана чрез общ модел на метаданни за полетата на филтрите. Всеки екран с търсене описва полетата си чрез надпис, описание, тип на входа и подсказващи текстове. Въз основа на тези метаданни преизползваемият визуализатор избира подходяща контрола, например текстово поле, падащ списък, числов интервал или интервал от дати. Въведените стойности не остават на ниво интерфейс, а се трансформират в типизирана заявка за търсене, която се използва от слоя за съхранение за изграждане на заявка към базата данни. Най-важното е, че тази функционалност не е написана за една конкретна таблица, а е демонстрирана върху два различни модела: \texttt{Videos} и \texttt{AnalysisRuns}.

По сходен начин е реализирана и преизползваемата функционалност за визуализиране на списъци от обекти. На преизползваемия list ViewModel се подава списък от DTO обекти, конфигурират се видимите колони и при избор на ред може да се върне целият избран обект, независимо кои свойства са били визуализирани. Визуализацията може да се изчиства и зарежда отново с нов списък. Същият механизъм е приложен върху \texttt{Videos} и \texttt{AnalysisRuns}, което показва, че таблицата и selection логиката не са обвързани с един конкретен модел.

\section{Ключови решения в кода}

Едно от най-важните архитектурни решения е превключването между SQLite и PostgreSQL. В проекта това не е реализирано чрез два различни слоя за съхранение, а чрез централизирана конфигурация в инфраструктурния слой. Така приложната логика остава една и съща, а конкретният доставчик се избира в стартовата конфигурация.

\lstinputlisting[
    firstline=16,
    lastline=83,
    caption={Откъс от \texttt{DependencyInjection.cs}, който избира SQLite или PostgreSQL},
    label={code:provider-switch}
]{../CVAnalysisHub.Infrastructure/DependencyInjection.cs}

Кодът от \cref{code:provider-switch} показва две важни идеи. Първо, низът за свързване се определя спрямо стойността на \texttt{Persistence:Provider}, а не е фиксиран в приложния слой. Второ, изборът между \texttt{UseSqlite} и \texttt{UseNpgsql} е концентриран на едно място. Именно това прави възможно изпълнението на едно и също приложение върху две СУБД без дублиране на бизнес логика.

Втората съществена част е слоят със споделени ViewModel-и. Следният откъс показва, че \texttt{VideoListViewModel} не е специфичен нито за Blazor, нито за Avalonia. Той комбинира споделени филтри, преизползваем списък от обекти и зареждане на данни през service contract, което му позволява да бъде използван от двата интерфейса.

\lstinputlisting[
    firstline=6,
    lastline=84,
    caption={Откъс от \texttt{VideoListViewModel.cs} със shared list и filter логика},
    label={code:video-list-vm}
]{../CVAnalysisHub.Application/Videos/VideoListViewModel.cs}

Откъсът от \cref{code:video-list-vm} показва как преизползваемата функционалност е капсулирана на ниво ViewModel, а не само на ниво визуален компонент. Списъкът от видеа, активните филтри, зареждането на резултатите и изборът на ред са организирани в една споделена приложна абстракция, която после се визуализира по различен начин от Blazor и Avalonia.

Преизползването не остава само в ViewModel слоя. \cref{code:object-list-table} показва и кратък реален пример от Blazor компонента за визуализация на списъци. Компонентът е generic, работи с \texttt{ObjectListViewModel<TItem>}, взима видимите колони от ViewModel-а и връща избрания ред като пълен обект, а не само като стойности от таблицата.

\lstinputlisting[
    firstline=1,
    lastline=60,
    caption={Откъс от \texttt{ObjectListTable.razor} като generic reusable list компонент},
    label={code:object-list-table}
]{../CVAnalysisHub.Web/Components/Shared/ObjectListTable.razor}

Следващият кратък откъс показва как задачите за анализ се отделят от потребителския интерфейс. \texttt{QueuedAnalysisBackgroundService} периодично търси чакащ анализ, създава scope за инфраструктурните услуги и делегира реалната обработка към \texttt{IAnalysisRunProcessor}. Така UI слоят създава задача, но не носи тежестта на самата обработка.

\lstinputlisting[
    firstline=8,
    lastline=35,
    caption={Откъс от \texttt{QueuedAnalysisBackgroundService.cs} за обработка на чакащи анализи},
    label={code:queued-worker}
]{../CVAnalysisHub.Infrastructure/AnalysisRuns/QueuedAnalysisBackgroundService.cs}

Последният ключов фрагмент е свързан с реалната последователност на инференцията. По-долу е показан съкратен откъс от модула, който интегрира YoloDotNet, извличането на кадри и генерирането на анотирано видео.

\lstinputlisting[
    firstline=23,
    lastline=98,
    caption={Откъс от \texttt{YoloDotNetAnalysisInferenceEngine.cs} с основния inference pipeline},
    label={code:yolo-engine}
]{../CVAnalysisHub.Infrastructure/AnalysisRuns/YoloDotNetAnalysisInferenceEngine.cs}

\cref{code:yolo-engine} показва, че анализът не се свежда до запис на фиктивни стойности в базата. Първо се валидира изходното видео и се създава временна работна директория. След това кадрите се извличат чрез \texttt{VideoProcessingService}, върху всеки кадър се извиква \texttt{RunObjectDetection}, резултатите се преобразуват в удобна за домейна структура, рисуват се ограничителни рамки и накрая се рендерира изходно анотирано видео. Така аналитичната стойност на проекта е свързана с изпълним цялостен процес, а не само с визуализация на предварително подготвени данни.

При Avalonia клиента визуализацията на медийното съдържание е реализирана чрез предварителни кадри и отваряне на оригиналното или обработеното видео със системния видео плейър. Така настолното приложение използва съществуващата \texttt{ffmpeg} последователност, без да въвежда тежка вградена зависимост за възпроизвеждане на видео.

\section{Трудности и преодолени проблеми}

В процеса на реализация възникнаха няколко трудности, които оказаха влияние върху стабилността и завършеността на приложението. Сред тях бяха синхронизацията между Web и Avalonia при различни доставчици и различно файлово хранилище, логиката по стартиране на настолния клиент и фоновите услуги, проблеми с темата и контраста в Avalonia под Linux, поведението на локалния HTTPS сертификат за разработка в браузъра и необходимостта ясно да се разграничат статусът на видеото и статусът на анализа в интерфейса.

Един от по-полезните практически изводи беше, че архитектурно правилният код сам по себе си не е достатъчен за завършено приложение. Необходимо е системата да бъде не само правилно структурирана, но и стабилна, разбираема и без подвеждащи състояния на интерфейса. Точно поради тази причина в по-късните етапи беше обърнато внимание и на удобството за работа с Avalonia клиента, така че той да бъде реално използваем интерфейс.
